% Part: incompleteness
% Chapter: incompleteness-provability
% Section: lob-thm

\documentclass[../../../include/open-logic-section]{subfiles}

\begin{document}

\olfileid[hi]{inc}{inp}{tar}

\olsection{सत्यता की अपरिभाष्यता}

\emph{परिभाष्यता} की धारणा अंकगणित की भाषा के लिए औपचारिक अर्थविज्ञान
होने पर निर्भर करती है। हमने अंकगणित की भाषा के सूत्रों और वाक्यों का एक
समुच्चय वर्णित किया है। ``अभिप्रेत व्याख्या'' ऐसे वाक्यों को प्राकृतिक
संख्याओं के बारे में दावे मानकर पढ़ती है, और ऐसा दावा सत्य या असत्य हो
सकता है। $\Struct{N}$ को वह !!{structure} मानें जिसका प्रांत $\Nat$ है और
जिसमें अंकगणित की भाषा के चिह्नों की मानक व्याख्या है। तब
$\Sat{N}{!A}$ का अर्थ है ``$!A$ मानक व्याख्या में सत्य है।''

\begin{defn}
प्राकृतिक संख्याओं का कोई संबंध $R(x_1,\dots,x_k)$,~$\Struct{N}$ में
\emph{परिभाष्य} तभी और केवल तभी है जब अंकगणित की भाषा में कोई सूत्र
$!A(x_1,\dots,x_k)$ ऐसा हो कि प्रत्येक $n_1,\dots,n_k$ के लिए
$R(n_1,\dots,n_k)$ तभी और केवल तभी जब
$\Sat{N}{!A(\num n_1,\dots,\num n_k)}$।
\end{defn}

भिन्न शब्दों में, कोई संबंध~$\Struct{N}$ में तभी और केवल तभी परिभाष्य है
जब वह सिद्धांत $\Th{TA}$ में निरूपणीय हो, जहाँ
$\Th{TA} = \Setabs{!A}{\Sat{N}{!A}}$ अंकगणित के सत्य वाक्यों का समुच्चय
है। (यदि यह आपको तुरंत स्पष्ट न हो, तो परिभाषाओं पर वापस जाकर जाँचें और
स्वयं को आश्वस्त करें कि ऐसा ही है।)

\begin{lem}
प्रत्येक संगणनीय संबंध~$\Struct{N}$ में परिभाष्य है।
\end{lem}

\begin{proof}
यह जाँचना सरल है कि $\Th{Q}$ में किसी संबंध का निरूपण करने वाला सूत्र
$\Struct{N}$ में उसी संबंध को परिभाषित करता है।
\end{proof}

अब पूछा जा सकता है कि क्या विलोम भी सत्य है? अर्थात् क्या~$\Struct{N}$ में
परिभाष्य प्रत्येक संबंध संगणनीय है? उत्तर नहीं है। उदाहरणतः:

\begin{lem}
विराम संबंध~$\Struct{N}$ में परिभाष्य है।
\end{lem}

\begin{proof}
याद करें कि क्लेनी सामान्य रूप प्रमेय कहता है कि प्रत्येक आंशिक संगणनीय
फलन~$f$ का कोई सूचकांक~$e$ ऐसा है कि सभी $x \in \Nat$ के लिए
$f(x) = U(\umin{s}{T(e,x,s)})$, जहाँ $U$ और $T$ आदिम पुनरावर्ती और इसलिए
सर्वत्र-परिभाषित हैं। अतः $f(x)$ तभी और केवल तभी परिभाषित है (अर्थात्
संगणना रुकती है) जब कोई~$s$ ऐसा हो कि $T(e,x,s)$ लागू हो।

अब $H$ को विराम संबंध मानें, अर्थात्
\[
H = \Setabs{\tuple{e,x}}{\lexists[s][T(e, x, s)]}.
\]
$!D_T$,~$\Struct{N}$ में $T$ को परिभाषित करे। तब
\[
H = \Setabs{\tuple{e,x}}{\Sat{N}{\lexists[s][!D_T(\num e, \num x, s)]}},
\]
इसलिए $\lexists[s][!D_T(z, x, s)]$,~$\Struct{N}$ में~$H$ को परिभाषित
करता है।
\end{proof}

\begin{prob}
दिखाएँ कि $Q(n) \defiff n \in \Setabs{\Gn{!A}}{\Th{Q} \Proves !A}$
अंकगणित में परिभाष्य है।
\end{prob}

स्वयं $\Th{TA}$ के बारे में क्या? क्या वह अंकगणित में परिभाष्य है? अर्थात्
क्या समुच्चय $\Setabs{\Gn{!A}}{\Sat{N}{!A}}$ अंकगणित में परिभाष्य है?
तार्स्की का प्रमेय इसका नकारात्मक उत्तर देता है।

\begin{thm}
\ollabel{thm:tarski}
अंकगणित के सत्य !!{sentence}s का समुच्चय अंकगणित में परिभाष्य नहीं है।
\end{thm}

\begin{proof}
मान लें $!D(x)$ उसे परिभाषित करता है, अर्थात् $\Sat{N}{!A}$ तभी और केवल
तभी जब $\Sat{N}{!D(\gn{!A})}$। नियत-बिंदु उपपत्ति से कोई सूत्र $!A$ ऐसा
है कि $\Th{Q} \Proves !A \liff \lnot !D(\gn{!A})$, और इसलिए
$\Sat{N}{!A \liff \lnot !D(\gn{!A})}$। किंतु तब $\Sat{N}{!A}$ तभी और
केवल तभी जब $\Sat{N}{\lnot !D(\gn{!A})}$, जो इस तथ्य के विरुद्ध है कि
$!D(y)$ को अंकगणित के सत्य कथनों का समुच्चय परिभाषित करना था।
\end{proof}

तार्स्की ने यह विश्लेषण सत्यता की अधिक सामान्य दार्शनिक धारणा पर लागू किया।
किसी भी भाषा $L$ के लिए तार्स्की ने तर्क दिया कि $L$ के लिए सत्यता की किसी
पर्याप्त धारणा को प्रत्येक वाक्य $X$ के लिए यह संतुष्ट करना होगा:
\begin{quote}
`$X$' तभी और केवल तभी सत्य है जब $X$।
\end{quote}
अंग्रेज़ी के लिए तार्स्की का प्रायः उद्धृत उदाहरण यह वाक्य है:
\begin{quote}
`बर्फ़ सफ़ेद है' तभी और केवल तभी सत्य है जब बर्फ़ सफ़ेद है।
\end{quote}
किंतु विकर्ण फलन का निरूपण करने योग्य प्रत्येक भाषा और प्रत्येक भाषिक
विधेय $T(x)$ के लिए हम ऐसा वाक्य $X$ बना सकते हैं जो ``$X$ तभी और केवल
तभी जब $T(\text{`$X$'})$ नहीं'' को संतुष्ट करता है। क्योंकि हम नहीं चाहते
कि सत्यता-विधेय कुछ वाक्यों को सत्य और असत्य दोनों घोषित करे, तार्स्की ने
निष्कर्ष निकाला कि किसी भाषा के सभी वाक्यों के लिए सत्यता-विधेय निर्दिष्ट
करने हेतु किसी न किसी ढंग से भाषा की सीमाओं के बाहर जाना आवश्यक है। दूसरे
शब्दों में, किसी भाषा का सत्यता-विधेय स्वयं उसी भाषा में परिभाषित नहीं किया
जा सकता।

\end{document}
